\documentclass[english,twoside, openright, hidelinks, 11 pt, class=report,crop=false]{standalone}
\input{../../preamb}
\input{../bm}

\begin{document}	
\opgt
\op{opgvekmidt}
Gitt punktene $ A=(m, n) $ og $ B=(s, t) $, og vektorene $ \vec{a}=[m, n] $ og $ \vec{b}=[s, t] $. Vis at midtpunktet til linjestykket $ AB $ er gitt ved uttrykket
\[ (0, 0) + \frac{1}{2}(\vec{u}+\vec{v}) \]

\op{vekfaktlen}
Gitt $ \vec{v}=[ca, cb] $. Vis at 
\[ |\vec{v}|=c\sqrt{a^2+b^2} \]

\op{opgvektenh} \vs
\abc{
\item Gitt en vektor $ \vec{v} $. Vis at lengden til vektoren $ \frac{\vec{v}}{|\vec{v}|} $ er lik 1.
\item Bestem uttrykket for vektoren som er parallell med vektoren $ [3, 4] $, og som har lengde 10.
}

\op{opgveklen}
Bestem lengden til hver av vektorene.
\alg{
\vec{a} &= [3, 4]\\
\vec{b} &= [-1, 7]\\
\vec{c} &= [-8, 6]\\
\vec{d} &= [4, -3]
}

\op{opgveknorm}
Undersøk om noen av vektorene fra \refops{opgveklen} står vinkelrett på hverandre.

\op{opgvekparl}
Undersøk om noen av vektorene fra \refops{opgveklen} er parallelle.

\newpage
\eksop{R1V22D1}{R1V22D1opg2}
For vektorene $ \vec{a} $ og $ \vec{b} $ er $ |\vec{a}|=2 $, $ \vec{b}=3 $ og $ \vec{a}\cdot\vec{b}=-3 $.\os

Vi lar $ \vec{u}=\vec{a}+\vec{b} $ og $ \vec{v}=\vec{a}-6\vec{b} $.
\abc{
\item Bestem lengden til $ \vec{u} $ og $ \vec{v} $.
\item Bestem vinkelen mellom $ \vec{u} $ og $ \vec{v} $.
}

\eksop{R1H25}{r1h25d1opg4}
Gitt punktene $A=(-2, 3)$ og $B=(3, 2)$ .
\abc{
\item Bestem lengden til $AB$.

\item Linja gjennom $A$ og $B$ skjærer $x$-aksen i punktet $C$. Bestem koordinatene til $C$.

\item $D=(2, t)$ der $t\in\mathbb{R}$. Bestem $t$ slik at $\angle ABD=90^\circ$. 
} 

\op{vekskalparl}
Gitt $ \vec{u}=[a,b] $ og $ \vec{v}=[c, d] $
Vis at hvis $\angle(\vec{u}, \vec{v})=0^\circ $, gir \eqref{skal2} at
\[ ad-bc=0 \]

\eksop{R1V22}{r1v22d1opg4}
Gitt tallet $t$ og punktene $A=(1, 2)$, $B=(-1, 5)$ og $C=(t, 4)$.
\abc{
\item Bestem $t$ slik at $\angle BAC=90^\circ$.
\item Bestem $t$ slik at $A$, $B$ og $C$ ligger på linje.
}

\eksop{R1V23D1}{R1V23D1opg3}
Gitt tre punkt $ A=(1, 3) $, $ B=(4, 0) $ og $ C=(9, 4) $.
\abc{
\item  Bruk vektorregning til å avgjøre om $ \angle CBA $ er mindre enn, lik eller større enn $ 90^\circ $.
}
Et punkt $ P $ ligger på linjen som går gjennom $ B $ og $ C $.
\abcs{1}{
\item Bruk vektorregning til å bestemme koordinatene til punktet P slik at $ AB\perp AP $.
}

\eksop{R1H23D1}{r1h23d1opg3}
I trekanten $ \triangle ABC $ er $ A= (-3, -1) $, $ B=(2, -2) $ og $ C=(5,2) $.
\abc{
\item Avgjør ved hjelp av vektorregning hvilken side i trekanten som er kortest. \item Avgjør ved hjelp av vektorrekning om noen av vinklane i trekanten er $ 90^\circ $. 
}

\eksop{R1V24}{r1v24d1opg4}
Gitt punktene $A=(3, 4)$, $B=(-1, -2)$ og $C=(3+t, 2t)$.
\abc{
\item Bestem $t$ slik at $A$, $B$ og $C$ ligger på linje.
\item Bestem $t$ slik at $\angle BCA=90^\circ $.	
}

\grubop{opgvekdiag}
Gitt et parallellogram $\square ABCD$. Vis at midpunktet til diagonalen $AC$ også er midtpunktet til diagonalen $BD$.

\grubop{r1v22d2opg2}
For vektorene $\vec{a}$ og $\vec{b}$ er $|\vec{a}|=2$, $|\vec{b}|=3$ og $\vec{a}\cdot\vec{b}=-3$. Gitt vektorene $\vec{u}=\vec{a}+\vec{b}$ og $\vec{v}=\vec{a}-6\vec{b}$.
\abc{
\item Bestem de respektive lengdene til $\vec{u}$ og $\vec{v}$.
\item Bestem vinkelen mellom $\vec{u}$ og $\vec{v}$.
}


\grubop{opgveklinperp}
Gitt en lineær funksjon $f$ med stigningstall $a$. Forklar hvorfor grafen til alle lineære funksjoner med stigningstall $-\frac{1}{a}$ vil stå vinkelrett på grafen til $f$.

\grubop{r1h22d1opg4} (R1h22)
Gitt punktene $A=(1, 1)$, $B=(9, 7)$ og $P=(5, 9)$.
\abc{
\item Vis at $\angle APB =90^\circ$.
\item En linje $l$ er parallell med $\vec{AB}$, og går gjennom punktet $P$. Det er også et annet punkt $Q$ på $l$ som er slik at $\angle AQB=90^\circ$. Bestem koordinatene til $Q$.
}

\grubop{r1h22d2opg3}
Gitt punktene $A=(0, 0)$, $B=(9, 1)$ og $C=(24, 10)$. En linje $l$ er gitt ved parameterfremstillingen
\begin{equation*}
	l(x)= \left\lbrace{
		\begin{array}{ll}
			x=12t \br
			y=5t
		\end{array}
	}\right.\qquad,\qquad t>0 
\end{equation*}
\abc{
\item Vis at $C$ ligger på $l$.
\item Finn koordinatene til punktet $D$ slik at $\angle ADB=120^\circ$.
\item Punktet $E$ ligger på $l$ slik at arealet til $\triangle ABE$ er 11. Finn koordinatene til $E$.
}

\grubop{r1h24d1opg5}
Gitt vektorene $\vec{u}=[3, -2]$, $\vec{v}=[4, 6]$, $\vec{w}=[2, -3]$ og $\vec{p}=[8, 12]$
\abc{
\item Avgjør om noen av vektorene er like lange.
\item Avgjør om noen av vektorene er ortogonale. 
\item Gitt vektoren $\vec{q}=[2a-3, 1+3b]$. Bestem $a$ og $b$ slik at $\vec{u}+2\vec{q}=[7, 5]$.
}

\grubop{r1h25d2opg5}
Gitt vektorene $\vec{a}$ og $\vec{b}$ hvor $|\vec{a}|=4$, $|\vec{b}|=2\sqrt{3}$, og $\angle (\vec{a}, \vec{b})=30^\circ$. For vektoren $\vec{p}$ har vi at $\vec{p}=\vec{a}+\vec{b}$.
\abc{
\item Finn lengden til $\vec{p}$.
\item For vektoren $\vec{q}$ har vi at $\vec{q}=t\vec{a}+\vec{b}$. Bestem $t$ slik at $\vec{p}$ og $\vec{q}$ er ortogonale.
}

\grubop{r1v25d1opg6} (R1V25)
Gitt punktene $O=(0, 0)$, $A=(-1, 2)$, og $B=(1, 1)$. 
\abc{
\item Finn lengden til $AB$.
\item Vektoren fra $O$ til $C=(-1, a)$ er parallell med $\vv{AB}$. Finn $a$.
\item Punktet $M$ er plassert slik at $\left|\vv{OM}\right|=\sqrt{10}$ og $\angle OBM=90^\circ$, og slik at avstanden mellom $M$ og $B$ er minst mulig. Finn $M$.
}

\grubop{r1h25d1opg4}
Gitt punktene $A=(-2, 3)$ og $B=(3, 2)$.
\abc{
\item Bestem lengden til $AB$.
\item Linjen gjennom $AB$ skjærer $x$-aksen i punktet $C$. Bestem koordinatene til $C$.
\item Gitt punktet $D=(2, t)$. Bestem $t$ slik at $\angle ABD=90^\circ$. 
}

\grubop{r1v23d2opg2}
Gitt punktet $A=(3, 2)$, tallet $t$ og vektorene $\vec{u}=[4, 3]$ og $\vec{v}=[2t, 5t]$. I parallellogremmet $\square ABCD$ er $\vv{AB}=\vec{u}$ og $\vv{AD}=\vec{v}$.
\abc{
\item Bestem koordinatene til $C$ og $D$ uttrykt ved $t$.
\item Bestem $t$ slik at skjæringspunktet mellom diagonalene i parallellogrammet er $P=(8, 11)$.
}


\end{document}